%! Polynomial Zeros and Factoring — Unit 3, Lesson 3.4 — AP Practice
% EXTRA CREDIT — optional, scored 1 to 5, exactly TWO pages, last in the packet (graded homework
% is DeltaMath). Page 1 is Section I, five multiple-choice items with four options (A)–(D); page
% 2 is Section II, one multi-part free-response question on a table. Its context (an open box
% folded from a 10-by-8 sheet) is used nowhere else in the lesson.
% Distractors are the lesson's errors on purpose: MC 1 divides by x and loses 0 (A, D) or
% forgets the square root (B); MC 2 stops early (A) or misreads the squares (C, D); MC 3 gets
% the sum-of-cubes signs wrong (B, D) or cubes the binomial (C); MC 4 drops the non-real pair
% (A), mis-signs it (B), or solves for x^2 instead of x (C); MC 5 divides by x^2 (A), ignores
% multiplicity (B), or forgets the square root (D).
%
% PAGE LOCKSTEP with the other half of the pair: the key marks the correct option in its
% LABEL (no text added to the line); every \gradedblank line is identical in both files. The
% explicit \newpage fixes the page break in both files. Every rubric box is drawn 2in tall.
\documentclass[10pt]{article}
\usepackage{precalculus-article}
\usepackage{precalculus-boxes}
\usepackage{grading}
\gradingassignment{precalc-u03-l04-ap}
\gradingcourse{precalculus}
\newcommand{\TallMath}[1]{$\displaystyle #1\rule[-1.4em]{0pt}{3.2em}$}

\begin{document}

\pageheader{Unit 3, Lesson 3.4}{AP Practice}
{\small\itshape Write each final answer in its box. Only what is inside a box is graded.}\par

\vspace{0.08in}

\begin{remindbox}
\small
\textbf{Extra credit, optional, scored 1 to 5.} These two pages are written in the style of the
AP Precalculus exam. Your graded homework is on DeltaMath; do it first. Every answer choice
below is a mistake someone really makes, so read all four before you choose.
\end{remindbox}

\vspace{0.08in}

\begin{headlinebox}{goldbg}
\small
\textbf{Section I --- Multiple Choice.} Write the letter of the single best answer in its
Letter box. No calculator. Work in the space under each question.
\end{headlinebox}

\vspace{0.04in}

\begin{enumerate}[leftmargin=1.9em, itemsep=24pt]

\item What are all the real solutions of $x^3 = 16x$?

      \smallskip
      \textbf{(A)} $-4$ and $4$ only \hspace{2em} \textbf{(B)} $0$ and $16$ \hspace{2em}
      \textbf{(C)} $-4$, $0$, and $4$ \hspace{2em} \textbf{(D)} $4$ only \hfill\textit{Letter:}~\gradedblank[id=ap1, type=choice, points=1, answer=C]{0.42in}{13pt}

\item Which of the following is the polynomial $x^4 - 1$ factored completely over the real
      numbers?

      \smallskip
      \begin{tabular}[b]{@{}p{5.6cm}p{5.6cm}@{}}
      \textbf{(A)} $(x^2 - 1)(x^2 + 1)$ & \textbf{(B)} $(x - 1)(x + 1)(x^2 + 1)$ \\[3pt]
      \textbf{(C)} $(x - 1)^2(x + 1)^2$ & \textbf{(D)} $(x - 1)(x + 1)(x^2 - 1)$
      \end{tabular}\hfill\textit{Letter:}~\gradedblank[id=ap2, type=choice, points=1, answer=B]{0.42in}{13pt}

\item Which of the following is equivalent to $x^3 + 64$?

      \smallskip
      \begin{tabular}[b]{@{}p{5.6cm}p{5.6cm}@{}}
      \textbf{(A)} $(x + 4)(x^2 - 4x + 16)$ & \textbf{(B)} $(x + 4)(x^2 + 4x + 16)$ \\[3pt]
      \textbf{(C)} $(x + 4)^3$ & \textbf{(D)} $(x - 4)(x^2 + 4x + 16)$
      \end{tabular}\hfill\textit{Letter:}~\gradedblank[id=ap3, type=choice, points=1, answer=A]{0.42in}{13pt}

\item What are all the solutions, real and non-real, of $x^4 - 2x^2 - 8 = 0$?

      \smallskip
      \begin{tabular}[b]{@{}p{5.6cm}p{5.6cm}@{}}
      \textbf{(A)} $2$ and $-2$ only & \textbf{(B)} $\pm 2$ and $\pm\sqrt{2}$ \\[3pt]
      \textbf{(C)} $4$ and $-2$ & \textbf{(D)} $\pm 2$ and $\pm i\sqrt{2}$
      \end{tabular}\hfill\textit{Letter:}~\gradedblank[id=ap4, type=choice, points=1, answer=D]{0.42in}{13pt}

\item Which of the following describes all the zeros of $p(x) = 2x^4 - 18x^2$?
      \begin{enumerate}[leftmargin=2.6em, itemsep=2pt]
        \item[(A)] $3$ and $-3$ only
        \item[(B)] $0$, $3$, and $-3$, each with multiplicity 1
        \item[(C)] $0$ with multiplicity 2, and $3$ and $-3$, each with multiplicity 1
        \item[(D)] $0$ with multiplicity 2, and $9$ with multiplicity 1 \hfill\textit{Letter:}~\gradedblank[id=ap5, type=choice, points=1, answer=C]{0.42in}{13pt}
      \end{enumerate}

\end{enumerate}

\newpage

\begin{headlinebox}{goldbg}
\small
\textbf{Section II --- Free Response.} Answer in context --- that earns credit on the AP exam.
Write each final answer in its box. Only what is inside a box is graded.
\end{headlinebox}

\vspace{0.04in}

\begin{enumerate}[leftmargin=1.9em, itemsep=2pt, topsep=2pt, start=6]

\item An open-top box is folded from a 10-inch by 8-inch sheet of cardboard by cutting a square
      $x$ inches on a side from each corner and folding up the sides. Its volume, in cubic
      inches, is $V(x) = x(10 - 2x)(8 - 2x) = 4x^3 - 36x^2 + 80x$. Some values of $V$ are
      shown in the table.

      \begin{center}
      \renewcommand{\arraystretch}{1.2}
      \begin{tabular}{|c|c|c|c|c|c|c|}
      \hline
      \rowcolor{lilac} $x$ (inches) & 0 & 1 & 2 & 3 & 4 & 5 \\ \hline
      $V(x)$ (cubic inches) & 0 & 48 & 48 & 24 & 0 & 0 \\ \hline
      \end{tabular}
      \end{center}

      \textbf{(a)} Use the factored form to solve $V(x) = 0$. For each solution, explain in
      context why that cut does not make a box.
      \rubricbox[2pt]{\gradedblank[id=ap6a, type=rubric, points=2,
        display={$x = 0$, $4$, $5$. A 0-inch cut leaves the sheet flat; a 4-inch cut uses up the 8-inch side ($8 - 2(4) = 0$); a 5-inch cut uses up the 10-inch side. A real box needs $0 < x < 4$.}]{\linewidth}{2in}}

      \textbf{(b)} Sam divides both sides of $4x^3 - 36x^2 + 80x = 0$ by $4x$, gets $x = 4$ or
      $x = 5$, and says $V$ has only two zeros. Is Sam correct? Explain what his method lost, and
      how the degree of $V$ shows something is missing.
      \rubricbox[2pt]{\gradedblank[id=ap6b, type=rubric, points=2,
        display={No. Dividing by $4x$ assumes $x \ne 0$, so it lost the zero $x = 0$. $V$ has degree 3, so it has three zeros: factor instead, $4x(x^2 - 9x + 20) = 4x(x - 4)(x - 5)$, giving $0$, $4$, $5$.}]{\linewidth}{2in}}

      \textbf{(c)} The table shows $V(1) = V(2) = 48$, and $V(x) = 48$ simplifies to
      $x^3 - 9x^2 + 20x - 12 = 0$. Verify by multiplying that this equals $(x - 1)(x - 2)(x - 6)$,
      give all three solutions, and explain why only two of them are possible cut sizes.
      \rubricbox[2pt]{\gradedblank[id=ap6c, type=rubric, points=2,
        display={$(x - 1)(x - 2) = x^2 - 3x + 2$; times $(x - 6)$: $x^3 - 6x^2 - 3x^2 + 18x + 2x - 12 = x^3 - 9x^2 + 20x - 12$. Solutions $x = 1, 2, 6$. A 6-inch cut is impossible: $8 - 2(6) < 0$, and a cut must satisfy $0 < x < 4$. So only 1 inch and 2 inches give a 48-cubic-inch box.}]{\linewidth}{2in}}

\end{enumerate}

\end{document}
